\documentclass[11pt]{article}
\usepackage[margin=1in]{geometry}
\usepackage[T1]{fontenc}
\usepackage[utf8]{inputenc}
\usepackage{lmodern}
\usepackage{graphicx}
\usepackage{booktabs}
\usepackage{longtable}
\usepackage{array}
\usepackage{hyperref}
\usepackage{xcolor}
\usepackage{float}
\usepackage{caption}
\hypersetup{colorlinks=true,linkcolor=blue,urlcolor=blue}

\title{ARC Solver Complexity Across Humans, LLMs, and Non-LLM Systems}
\author{}
\date{April 7, 2026}

\begin{document}
\maketitle

\section*{What This Paper Is}

This paper consolidates the full analysis thread in this folder into one place. It starts from the approved ARC Python solvers, explains how they were fetched and validated, summarizes the complexity panel that was computed over those solvers, and then traces how those complexity measures relate to humans, LLMs, thinking-vs-standard model differences, and the non-LLM systems in the repo.

The central question is:

\begin{quote}
When a task requires a more structurally complex correct Python solver, does that show up as greater task difficulty for humans, LLMs, or both?
\end{quote}

The short answer is:

\begin{itemize}
  \item yes for LLMs, strongly
  \item yes for humans, but much more weakly
  \item yes for non-LLM systems directionally, but not yet with enough power to be confident
\end{itemize}

The strongest significant contrast in the whole project is:

\begin{quote}
Structural solver complexity is more strongly associated with LLM difficulty than with human difficulty on the matched overlap tasks.
\end{quote}

\section*{Executive Summary}

\begin{itemize}
  \item We fetched \texttt{511} ARC \texttt{solution.py} files and validated them against the official ARC task data.
  \item Only \texttt{127 / 511} passed validation.
  \item The site's \texttt{approved} status explained almost all of that gap: \texttt{120 / 120} approved programs passed, while nearly all \texttt{submitted}, \texttt{attempted}, and \texttt{skipped} programs failed.
  \item We therefore restricted the main analysis to \texttt{120} validated approved Python solvers.
  \item Across those solvers, structural complexity measures like \texttt{ast\_node\_count}, \texttt{token\_count}, \texttt{cyclomatic\_complexity}, and \texttt{halstead\_volume} clustered tightly together.
  \item The best single solver-complexity proxy for latent difficulty was \texttt{ast\_node\_count} ($r = 0.666$).
  \item A conservative supervised composite improved that only modestly (PCR-5, leave-one-out $r = 0.691$).
  \item On the expanded LLM analysis, \texttt{log1p(cyclomatic\_complexity)} reached $r = 0.707$ against simple LLM logit difficulty.
  \item Humans and LLMs share a real common axis on the approved overlap tasks:
  \begin{itemize}
    \item human difficulty vs LLM difficulty: $r = 0.531$, $p = 0.0268$, $q = 0.0366$
    \item human solve rate vs LLM pass rate: $r = 0.541$, $p = 0.0187$, $q = 0.0280$
  \end{itemize}
  \item The strongest significant difference claim is:
  \begin{itemize}
    \item cyclomatic complexity is more strongly associated with LLM difficulty than with human difficulty: $\Delta r = 0.441$, $p = 0.0378$, $q = 0.0472$
  \end{itemize}
  \item Residual analyses sharpen that story:
  \begin{itemize}
    \item residual LLM difficulty after removing human difficulty still tracks cyclomatic complexity: $r = 0.603$, $p = 0.0117$, $q = 0.0194$
    \item residual human difficulty after removing LLM difficulty tracks human duration only suggestively: $r = 0.470$, $p = 0.0505$, $q = 0.0583$
  \end{itemize}
  \item Human pair-level data show a different signature from LLM difficulty:
  \begin{itemize}
    \item human pair difficulty tracks human duration: $r = 0.391$, $p = 1.67\times 10^{-4}$, $q = 3.57\times 10^{-4}$
    \item raw board size alone is weak for human difficulty: $r = -0.096$, $p = 0.311$
    \item within-task human difficulty heterogeneity is substantial: mean range $0.719$, max $2.498$
  \end{itemize}
  \item The negative \texttt{thinking\_advantage} result did not survive the full scientific audit cleanly. It looked strong under the legacy grouping, but weakened materially after verified relabeling and floor-sensitive checks.
  \item The first-pass non-LLM results are interesting but underpowered. Their complexity correlations are directionally between human and LLM correlations on the shared \texttt{17} ARC-2 tasks, but the intermediate position is not yet statistically distinguishable from either side.
\end{itemize}

\section*{Data Assembly and Validation}

\subsection*{Task IDs and local ARC data}

The local ARC mirror was checked against the public ARC-AGI-1 and ARC-AGI-2 repositories.

\begin{itemize}
  \item ARC-AGI-1 matched the official \texttt{400} train and \texttt{400} eval task IDs exactly.
  \item ARC-AGI-2 matched the official \texttt{1000} train and \texttt{120} eval task IDs exactly.
  \item ARC-AGI-1 train and eval plus ARC-AGI-2 train reconstructed cleanly against the official public task content.
  \item ARC-AGI-2 eval had \texttt{6} task-content differences versus the current public repo: \texttt{4a21e3da}, \texttt{abc82100}, \texttt{b6f77b65}, \texttt{d8e07eb2}, \texttt{f560132c}, and \texttt{faa9f03d}.
\end{itemize}

This mattered for provenance, but not for the basic task-ID fetch pipeline.

\subsection*{Solver fetch and validation}

We discovered \texttt{1,147} unique ARC task IDs across ARC-1 train/eval and ARC-2 train/eval and used the direct site endpoint \texttt{https://arc.huikang.dev/solutions/<task\_id>/solution.py}.

Fetch outcome:

\begin{itemize}
  \item \texttt{511} unique \texttt{solution.py} files on disk
  \item \texttt{636} tasks with no stored \texttt{solution.py}
\end{itemize}

Validation outcome on the official task data:

\begin{itemize}
  \item passed: \texttt{127}
  \item wrong answer: \texttt{380}
  \item crash: \texttt{4}
  \item timeout: \texttt{0}
\end{itemize}

The site status labels explained the pattern almost perfectly:

\begin{itemize}
  \item \texttt{approved}: \texttt{120 / 120} passed
  \item \texttt{submitted}: \texttt{1 / 236} passed
  \item \texttt{attempted}: \texttt{5 / 138} passed, \texttt{4} crashed
  \item \texttt{skipped}: \texttt{1 / 17} passed
  \item \texttt{correct}: \texttt{629} tasks had no stored \texttt{solution.py}
\end{itemize}

That gave a clean methodological decision: use only the \texttt{approved} solutions for the complexity study.

\subsection*{Approved package}

The final approved-only package contains \texttt{120} unique task IDs. By dataset membership:

\begin{itemize}
  \item ARC-AGI 1 train: \texttt{57}
  \item ARC-AGI 1 eval: \texttt{39}
  \item ARC-AGI 2 train: \texttt{100}
  \item ARC-AGI 2 eval: \texttt{20}
\end{itemize}

\section*{Complexity Measures}

\subsection*{Static, structural, and dynamic metrics}

The solver complexity panel included:

\begin{itemize}
  \item text/size proxies: \texttt{nonblank\_lines}, \texttt{token\_count}, \texttt{gzip\_bytes}
  \item structural proxies: \texttt{ast\_node\_count}, \texttt{cyclomatic\_complexity}, \texttt{branch\_node\_count}, \texttt{max\_nesting\_depth}
  \item description-length style proxies: \texttt{halstead\_volume}, \texttt{halstead\_effort}
  \item dynamic/runtime proxies: \texttt{opcode\_count\_dynamic}, \texttt{elapsed\_ms\_total}, \texttt{elapsed\_ms\_per\_test}, and normalized runtime measures
\end{itemize}

On the approved solver set:

\begin{itemize}
  \item nonblank LOC median: \texttt{55}
  \item nonblank LOC mean: \texttt{67.8}
  \item cyclomatic complexity median: \texttt{17.5}
  \item cyclomatic complexity mean: \texttt{23.9}
  \item Halstead volume median: \texttt{1623}
  \item maximum nesting depth median: \texttt{4}
\end{itemize}

The important result was not the raw medians, but that static structural metrics clustered tightly while runtime measures formed a weaker, partially separate axis.

\subsection*{PCA and composite models}

PCA over the main complexity panel found:

\begin{itemize}
  \item PC1: \texttt{57.4\%} of variance
  \item PC2: \texttt{17.5\%}
  \item PC3: \texttt{9.8\%}
\end{itemize}

Interpretation:

\begin{itemize}
  \item PC1 reads as overall solver size and structural density.
  \item PC2 reads as runtime intensity and execution burden.
  \item PC3 reads more like grid volume and memory footprint.
\end{itemize}

The best single versus composite predictors were:

\begin{itemize}
  \item best single metric: \texttt{ast\_node\_count}, $r = 0.666$
  \item unsupervised PC1: $r = 0.629$
  \item best conservative composite: PCR-5, leave-one-out $r = 0.691$
\end{itemize}

\begin{figure}[H]
  \centering
  \includegraphics[width=0.9\linewidth]{chart_pca_overview.png}
  \caption{PCA view of the complexity panel. A dominant structural component coexists with secondary runtime and grid-size components.}
\end{figure}

\begin{figure}[H]
  \centering
  \includegraphics[width=0.9\linewidth]{chart_prediction_comparison.png}
  \caption{Best single metric versus a conservative composite predictor. The composite helps, but only modestly.}
\end{figure}

\section*{Human and LLM Psychometric Data}

\subsection*{LLM-side overlap}

The main LLM psychometric join covered \texttt{56} approved eval rows:

\begin{itemize}
  \item \texttt{38} ARC-1 eval
  \item \texttt{18} ARC-2 eval
\end{itemize}

LLM difficulty was summarized in several ways: earlier shared-model latent difficulty, pooled all-model Rasch difficulty, simple pass-rate and logit-difficulty summaries, and 2PL difficulty/discrimination.

These summaries were extremely stable:

\begin{itemize}
  \item previous latent difficulty vs pooled Rasch: $r = 0.991$
  \item pooled Rasch vs simple logit difficulty: $r = 0.970$
\end{itemize}

\subsection*{Human-side overlap}

The human task-level overlap came from the public evaluation table and is much narrower. The key limitation is that the human task table is almost entirely ARC-2 public eval. Only \texttt{6} human tasks belong to both ARC-1 eval and ARC-2 eval, and only \texttt{2} approved ARC-1 eval tasks overlap with the approved human-linked task set.

So the fully matched human + LLM + non-LLM + solver-complexity comparison is effectively an ARC-2 comparison on \texttt{17} independent tasks. That is a data-overlap limit, not a coding bug.

\section*{Main LLM Difficulty Results}

The strongest LLM-side result is that structural solver complexity tracks LLM difficulty strongly.

\begin{itemize}
  \item \texttt{log1p(cyclomatic\_complexity)} vs LLM logit difficulty: $r = 0.707$
  \item \texttt{ast\_node\_count} vs pooled Rasch difficulty: $r = 0.653$
  \item \texttt{log1p(cyclomatic\_complexity)} vs previous shared latent difficulty: $r = 0.691$
\end{itemize}

Structural code measures consistently outperformed raw runtime burden for LLM difficulty. In ARC, LLM difficulty appears to track something closer to structural rule complexity and compact solver description length than brute execution effort.

\begin{figure}[H]
  \centering
  \includegraphics[width=0.92\linewidth]{chart_llm_complexity_heatmap.png}
  \caption{LLM difficulty and related psychometric outcomes against the full complexity panel. The difficulty-family signals are dominated by structural metrics.}
\end{figure}

\section*{Human vs LLM: Shared Axis and Significant Difference}

\subsection*{Shared axis}

Humans and LLMs are not living on unrelated difficulty scales.

\begin{center}
\small
\begin{tabular}{p{3.0in}cccc}
\toprule
Claim & Estimate & 95\% CI & p & q \\
\midrule
Human difficulty aligns with LLM difficulty (S1) & $r=0.531$ & $[0.178,0.780]$ & 0.0268 & 0.0366 \\
Human solve rate aligns with LLM pass rate (S2) & $r=0.541$ & $[0.226,0.781]$ & 0.0187 & 0.0280 \\
Previous latent vs pooled Rasch (S3) & $r=0.991$ & $[0.984,0.997]$ & $1.67\times10^{-4}$ & $3.57\times10^{-4}$ \\
Pooled Rasch vs simple logit difficulty (S4) & $r=0.970$ & $[0.953,0.981]$ & $1.67\times10^{-4}$ & $3.57\times10^{-4}$ \\
\bottomrule
\end{tabular}
\end{center}

\subsection*{The key significant difference}

The most important result in the entire project is:

\begin{itemize}
  \item D1: cyclomatic complexity is more strongly associated with LLM difficulty than with human difficulty
  \item $\Delta r = 0.441$
  \item 95\% CI $[0.024, 0.909]$
  \item $p = 0.0378$
  \item $q = 0.0472$
\end{itemize}

The raw correlations behind that difference are:

\begin{itemize}
  \item cyclomatic vs human difficulty: $r = 0.150$
  \item cyclomatic vs LLM difficulty: $r = 0.591$
\end{itemize}

That is the clearest significant evidence that human difficulty and LLM difficulty are not just noisy copies of the same thing.

\subsection*{Residual split}

Residual analyses sharpen the same story:

\begin{itemize}
  \item D4: residual LLM difficulty after removing human difficulty still tracks cyclomatic complexity
  \item $r = 0.603$
  \item 95\% CI $[0.284, 0.816]$
  \item $p = 0.0117$
  \item $q = 0.0194$
\end{itemize}

But the analogous human-side result is only suggestive:

\begin{itemize}
  \item D3: residual human difficulty vs mean human duration
  \item $r = 0.470$
  \item 95\% CI $[0.067, 0.766]$
  \item $p = 0.0505$
  \item $q = 0.0583$
\end{itemize}

So the pattern is:

\begin{itemize}
  \item shared axis exists
  \item LLM-specific residuals still look structural
  \item human-specific residuals look more time/search-like, but the task-level overlap is small
\end{itemize}

\begin{figure}[H]
  \centering
  \includegraphics[width=0.95\linewidth]{chart_synthesis_throughline.png}
  \caption{Main through-line figure. Humans and LLMs share a common axis, but the residual structure splits into a more structural LLM side and a more time/search-like human side.}
\end{figure}

\section*{Human-Side Findings Beyond the Overlap Table}

The pair-level human data are where the human story becomes clearer.

\begin{itemize}
  \item H1: human pair difficulty tracks mean human duration: $r = 0.391$, $p = 1.67\times10^{-4}$, $q = 3.57\times10^{-4}$
  \item H2: raw board size alone is weak for human difficulty: $r = -0.096$, $p = 0.311$
  \item H3: duration is more informative than raw board size: $\Delta r = 0.487$, $p = 1.25\times10^{-4}$, $q = 3.57\times10^{-4}$
  \item H4: the human-over-LLM gap shrinks when public-eval tasks expose more test pairs: $r = -0.366$, $p = 3.33\times10^{-4}$, $q = 6.25\times10^{-4}$
  \item H5: task identity explains a large share of pair-level human difficulty variation on multi-pair tasks: $R^2 = 0.749$, $p = 7.93\times10^{-5}$, $q = 3.57\times10^{-4}$
\end{itemize}

Descriptive heterogeneity results:

\begin{itemize}
  \item \texttt{44} public-eval tasks had multiple test pairs
  \item mean within-task difficulty range: $0.719$
  \item max within-task difficulty range: $2.498$
\end{itemize}

This matters because a single solver file is a task-level object, while human difficulty can vary meaningfully across test pairs within the same task. That gives a plausible reason why human difficulty should not be expected to line up with solver structure as strongly as LLM difficulty does.

\begin{figure}[H]
  \centering
  \includegraphics[width=0.95\linewidth]{chart_synthesis_human_pair_level.png}
  \caption{Human-side pair-level structure. Human difficulty is more time/search-like than board-size-like, and within-task heterogeneity is substantial.}
\end{figure}

\section*{Thinking Advantage: Initial Result, Audit, and Downgrade}

The original definition was \texttt{thinking\_advantage = pass\_rate\_thinking - pass\_rate\_standard}, with a logit-gap analog.

Under the original heuristic model grouping, the result looked strong:

\begin{itemize}
  \item T1: thinking advantage declines with LLM difficulty: $r = -0.492$, $p = 1.67\times10^{-4}$, $q = 3.57\times10^{-4}$
  \item T2: grouped binomial GLM interaction is negative: coefficient $-0.718$, $p = 5.70\times10^{-6}$, $q = 8.55\times10^{-5}$
\end{itemize}

We then did two things:

\begin{enumerate}
  \item audited ambiguous thinking-model labels against provider documentation
  \item checked whether the result was mainly a floor artifact
\end{enumerate}

Ambiguous models included \texttt{gemini-3-pro-preview}, \texttt{gpt-5-pro-2025-10-06}, \texttt{gpt-5-2-pro-2025-12-11-high}, \texttt{gpt-5-2-pro-2025-12-11-medium}, and \texttt{QwQ-32B-Fireworks}.

Under the verified grouping:

\begin{itemize}
  \item standard-group zero-success items: \texttt{35 / 56}
  \item thinking-group zero-success items: \texttt{0 / 56}
  \item both-zero items: \texttt{0 / 56}
\end{itemize}

The result then weakened sharply on floor-resistant subsets:

\begin{itemize}
  \item all rows: raw-gap $r = -0.378$
  \item standard-nonzero subset: raw-gap $r = 0.280$
  \item fully interior subset: raw-gap $r = -0.009$
\end{itemize}

So the scientifically responsible read is:

\begin{itemize}
  \item the original result was not pure nonsense
  \item but it is not robust enough after label verification and floor checks to treat as a stable substantive conclusion
\end{itemize}

\begin{figure}[H]
  \centering
  \includegraphics[width=0.95\linewidth]{chart_synthesis_thinking_audit.png}
  \caption{Thinking-advantage audit. The legacy negative result weakens substantially after verified relabeling and floor-sensitive subsetting.}
\end{figure}

\section*{Non-LLM First Pass}

\subsection*{Data used}

The non-LLM first pass used:

\begin{itemize}
  \item \texttt{38} approved ARC-1 eval tasks
  \item \texttt{18} approved ARC-2 eval tasks
  \item \texttt{17} shared ARC-2 tasks with human, LLM, and non-LLM outcomes
\end{itemize}

The system profiles included ARC-1 VARC and CompressARC variants and ARC-2 TRM, VARC, and related non-LLM profiles.

\subsection*{Main finding}

The non-LLM results were directionally coherent but underpowered.

\begin{itemize}
  \item ARC-2 non-LLM difficulty vs \texttt{gzip\_bytes}: $r = 0.504$, $q = 0.217$
  \item ARC-1 non-LLM difficulty vs \texttt{gzip\_bytes}: $r = 0.366$, $q = 0.191$
  \item pooled within-dataset-standardized non-LLM difficulty vs \texttt{gzip\_bytes}: $r = 0.378$, $q = 0.062$
\end{itemize}

None of the pre-specified non-LLM key tests N1 through N9 survived BH-FDR correction.

\subsection*{The interesting in-between pattern}

On the shared \texttt{17} ARC-2 tasks:

\begin{center}
\small
\begin{tabular}{p{2.4in}ccc}
\toprule
Metric and outcome & $r$ & $p$ & 95\% CI \\
\midrule
Cyclomatic vs human difficulty & 0.150 & 0.565 & $[-0.356, 0.588]$ \\
Cyclomatic vs non-LLM difficulty & 0.356 & 0.161 & $[-0.151, 0.714]$ \\
Cyclomatic vs LLM difficulty & 0.591 & 0.0125 & $[0.154, 0.834]$ \\
PC1 vs human difficulty & 0.135 & 0.604 & $[-0.369, 0.578]$ \\
PC1 vs non-LLM difficulty & 0.388 & 0.124 & $[-0.114, 0.732]$ \\
PC1 vs LLM difficulty & 0.535 & 0.0268 & $[0.074, 0.808]$ \\
\bottomrule
\end{tabular}
\end{center}

So non-LLM systems look visually intermediate. But the paired difference tests say we should be careful:

\begin{itemize}
  \item for cyclomatic complexity:
  \begin{itemize}
    \item LLM $>$ human: supported, $p = 0.039$
    \item LLM $>$ non-LLM: not supported, $p = 0.298$
    \item non-LLM $>$ human: not supported, $p = 0.569$
  \end{itemize}
  \item for complexity PC1:
  \begin{itemize}
    \item no pairwise difference was significant
  \end{itemize}
\end{itemize}

\subsection*{Power reality}

This is mostly a power problem. With \texttt{n = 17}, a two-sided \texttt{.05} correlation study needs about $|r| \approx 0.63$ for \texttt{80\%} power.

Approximate sample sizes needed for \texttt{80\%} power:

\begin{itemize}
  \item $r = 0.30$: $n \approx 85$
  \item $r = 0.35$: $n \approx 62$
  \item $r = 0.40$: $n \approx 47$
  \item $r = 0.50$: $n \approx 30$
\end{itemize}

So the non-LLM ``middle'' pattern is plausible, but we cannot yet separate it cleanly from either side.

\begin{figure}[H]
  \centering
  \includegraphics[width=0.9\linewidth]{chart_non_llm_complexity_comparison.png}
  \caption{Human vs LLM vs non-LLM complexity comparison. The non-LLM systems appear intermediate, but the overlap is too small for a strong three-way ordering claim.}
\end{figure}

\section*{Statistical Framework}

\subsection*{Estimators used}

\begin{itemize}
  \item Pearson correlations
  \item bootstrap 95\% confidence intervals
  \item permutation p-values for correlation tests
  \item BH-FDR correction for the main named claim families
  \item bootstrap differences of correlations for matched human-vs-LLM and related contrasts
  \item grouped binomial GLMs for thinking-vs-standard analyses
  \item OLS residualization for shared-vs-specific analyses
\end{itemize}

\subsection*{Null hypotheses}

\begin{itemize}
  \item shared-axis null: two task orderings are unrelated
  \item difference-of-correlation null: a predictor is equally associated with two outcomes
  \item residual null: after removing the shared axis, the remaining residual is unrelated to the tested predictor
  \item pair-level human null: the pair-level human measure is unrelated to the tested pair feature
  \item thinking-gap null: thinking and standard model groups have the same difficulty slope
\end{itemize}

\section*{Through-Line}

The results do not feel random. They line up into a coherent story:

\begin{enumerate}
  \item The approved solver set is real and clean.
  \item Structural complexity dominates the solver complexity panel.
  \item LLM difficulty is extremely stable across different psychometric constructions.
  \item Human and LLM difficulty share a real common axis.
  \item The cleanest significant difference is that structural solver complexity is more LLM-linked than human-linked.
  \item Residual LLM difficulty still looks structural.
  \item Human difficulty looks more duration/search-like and more pair-heterogeneous.
  \item The non-LLM systems seem intermediate, but current overlap is too small for a decisive claim.
  \item The one especially flashy result, thinking advantage, became weaker when audited properly.
\end{enumerate}

That pattern is exactly what a sane pipeline should look like: some results strengthen, some stay suggestive, and one attractive result gets downgraded.

\begin{figure}[H]
  \centering
  \includegraphics[width=0.92\linewidth]{chart_synthesis_stats_forest.png}
  \caption{Key claim estimates with uncertainty. The human-vs-LLM structural split and the human pair-level timing results are the cleanest supported findings.}
\end{figure}

\section*{Full Hypothesis Checklist}

\subsection*{Shared human-LLM alignment}

\begin{itemize}
  \item S1: Human and LLM task difficulty are positively aligned on the approved ARC-2 overlap. Result: supported. $r = 0.531$, $p = 0.0268$, $q = 0.0366$.
  \item S2: Human solve rate aligns with average-model pass rate on the same overlap tasks. Result: supported. $r = 0.541$, $p = 0.0187$, $q = 0.0280$.
  \item S3: The earlier shared-model latent scale and the new pooled Rasch scale are effectively the same LLM difficulty axis. Result: strongly supported. $r = 0.991$, $p = 1.67\times10^{-4}$, $q = 3.57\times10^{-4}$.
  \item S4: Pooled Rasch difficulty and simple LLM logit difficulty are almost identical. Result: strongly supported. $r = 0.970$, $p = 1.67\times10^{-4}$, $q = 3.57\times10^{-4}$.
\end{itemize}

\subsection*{Human vs LLM difference claims}

\begin{itemize}
  \item D1: Cyclomatic complexity is more strongly associated with LLM difficulty than with human difficulty. Result: supported. $\Delta r = 0.441$, $p = 0.0378$, $q = 0.0472$.
  \item D2: Human duration is more strongly associated with human difficulty than with LLM difficulty. Result: suggestive, not FDR-supported. $\Delta r = 0.374$, $p = 0.0820$, $q = 0.0879$.
  \item D3: Residual human difficulty still tracks human duration after removing shared LLM difficulty. Result: suggestive, not FDR-supported. $r = 0.470$, $p = 0.0505$, $q = 0.0583$.
  \item D4: Residual LLM difficulty still tracks solver structure after removing shared human difficulty. Result: supported. $r = 0.603$, $p = 0.0117$, $q = 0.0194$.
\end{itemize}

\subsection*{Human pair-level claims}

\begin{itemize}
  \item H1: Human pair difficulty tracks human time cost. Result: supported.
  \item H2: Raw board size alone is weak for human difficulty. Result: not supported as a positive predictor, which is the point.
  \item H3: Human duration is more informative than raw board size for human difficulty. Result: supported.
  \item H4: Human-over-LLM advantage shrinks when tasks expose more test pairs. Result: supported.
  \item H5: Task identity explains a large share of pair-level human difficulty variance on multi-pair tasks. Result: supported.
\end{itemize}

\subsection*{Thinking-advantage claims}

\begin{itemize}
  \item T1: Thinking advantage declines as approved-item difficulty rises under the legacy label schema. Result: statistically supported under the legacy schema, but downgraded after label audit and floor-sensitive checks.
  \item T2: Thinking-vs-standard success probability has a negative difficulty interaction under the legacy label schema. Result: statistically supported under the legacy schema, but not treated as a final substantive claim after the verified-label audit.
\end{itemize}

\subsection*{Non-LLM key tests}

\begin{itemize}
  \item N1: ARC-2 non-LLM difficulty aligns with ARC-2 LLM difficulty. Result: not supported.
  \item N2: ARC-2 non-LLM difficulty aligns with human difficulty. Result: not supported.
  \item N3: Human difficulty is more strongly aligned with LLM difficulty than with non-LLM difficulty. Result: not supported.
  \item N4: ARC-2 non-LLM difficulty positively tracks structural solver complexity. Result: directionally positive but not supported after correction.
  \item N5: Structural solver complexity is more strongly associated with LLM difficulty than with non-LLM difficulty. Result: not supported.
  \item N6: After controlling LLM difficulty, non-LLM difficulty still retains structural-complexity signal. Result: not supported.
  \item N7: For non-LLM difficulty, structural complexity is more informative than runtime intensity. Result: not supported.
  \item N8: ARC-1 non-LLM difficulty positively tracks structural solver complexity. Result: directionally positive but not supported after correction.
  \item N9: CompressARC search-depth difficulty positively tracks structural solver complexity. Result: not supported.
\end{itemize}

\subsection*{Supplemental ARC-2 overlap-only checks}

\begin{itemize}
  \item On the \texttt{17} matched ARC-2 tasks, LLM correlations with structure are the only ones that are individually significant in the human/LLM/non-LLM side-by-side comparison.
  \item The ``non-LLM sits in the middle'' pattern is descriptively real-looking but underpowered.
\end{itemize}

\section*{Limitations}

\begin{itemize}
  \item The approved solver set is only \texttt{120} tasks.
  \item Human task-level overlap with the approved set is only \texttt{17} independent tasks.
  \item Human public-eval data are mostly ARC-2, so ARC-1 does not add much independent human-overlap power.
  \item Many complexity metrics are highly collinear, so multiple large correlations do not imply many independent discoveries.
  \item Observed solver size is not minimal description length; it is the size of one approved solver implementation.
  \item The thinking analysis depends on model grouping and bounded pass-rate scales, so it needed more auditing than the other analyses.
\end{itemize}

\section*{Bottom Line}

The strongest defensible statement from this project is:

\begin{quote}
In approved ARC solvers, structural program complexity is a strong proxy for LLM task difficulty. Humans and LLMs share part of the same difficulty axis, but structural solver complexity is significantly more LLM-linked than human-linked, while human-specific difficulty looks more time/search-like and more pair-heterogeneous.
\end{quote}

The best next step, if more certainty is needed, is not to add more complexity metrics. It is to increase the matched human-overlap sample or find a second benchmark with the same low-boilerplate solver structure.

\vspace{0.5em}
\noindent Main files used in this paper:
\begin{itemize}
  \item \texttt{complexity\_master\_hypotheses.csv}
  \item \texttt{synthesis\_writeup.md}
  \item \texttt{statistical\_hypothesis\_report.md}
  \item \texttt{human\_llm\_difference\_report.md}
  \item \texttt{human\_additional\_findings.md}
  \item \texttt{approved\_llm\_complexity\_report.md}
  \item \texttt{non\_llm\_complexity\_addendum.md}
  \item \texttt{non\_llm\_arc2\_overlap\_significance.md}
\end{itemize}

\end{document}
